\documentclass[sigconf,review]{acmart}
\usepackage{graphicx} % Required for inserting images
\usepackage{booktabs}
\usepackage{tabularx}
\usepackage{multirow}
\usepackage{listings}
\usepackage{mdframed}
\usepackage{xcolor}

\definecolor{boxgray}{RGB}{245,245,245}
\definecolor{boxborder}{RGB}{180,180,180}
\lstset{basicstyle=\ttfamily\scriptsize, breaklines=true,
        frame=none, backgroundcolor=\color{boxgray}}

\newcommand{\note}[1]{\textcolor{magenta}{{{[#1]}}}}
\newcommand{\he}[1]{\textcolor{purple}{{{\bfseries [[hannah: #1]]}}}}

\title{SE4AI - Requirements Level Explanations for Different Stakeholders}
\author{Hannah Estes}
\author{Mohsen EsfandyariDoulabi}
\author{Saba Khanmohammadi}
\date{March 2026}

\begin{document}

\maketitle

\section{Introduction}
Software engineers routinely face decisions that require balancing competing objectives, such as faster delivery versus lower cost, security versus usability, and accuracy versus interpretability. Multi-objective optimization offers principled methods for navigating these trade offs, and repositories such as MOOT~\cite{menzies2025moot} now provide over 120 real world SE optimization tasks spanning software configuration, cloud tuning, project health, and process modeling. The algorithmic machinery exists; the interpretation problem does not. 
When an optimizer produces a Pareto front or a decision tree over configuration parameters, a developer or project manager sees numeric trade offs between objective values, throughput versus latency, effort versus defect risk. What they do not see is what those trade-offs mean at the level of requirements: which stakeholder concerns are being satisfied, which are being compromised, and why a particular configuration represents a defensible decision. This interpretation gap between optimization outputs and requirement-level reasoning is a persistent barrier to adoption. As Menzies et al. observe, industrial practitioners continue to deliver sub-optimal products in part because they lack the tools to explore and understand trade offs.

Recent work in explainable AI (XAI) has made progress on translating model outputs into natural language, and large language models (LLMs) have shown promise as verbalization engines for decision-tree explanations ~\cite{serafim2024exploring}. Separately, the requirements engineering literature has long recognized that stakeholders reason about trade-offs differently depending on their role, priorities, and technical background ~\cite{hujainah2018software}. However, no prior work combines these threads: generating stakeholder-tailored, requirement-level explanations of multi-objective optimization outputs, and evaluating whether such explanations actually improve trade off reasoning.     
This paper addresses that gap. We present a framework that maps multi-objective optimization outputs from MOOT tasks to requirement-level explanations via a structured schema and LLM-based verbalization. We evaluate explanation quality along three dimensions: faithfulness, clarity, and consistency, and assess whether tailoring explanations to distinct stakeholder profiles changes how trade offs are understood and acted upon. Specifically, we ask:

\begin{itemize}
    \item \textbf{RQ1:} How can multi-objective optimization outputs be systematically translated into requirement-level explanations?
    \item \textbf{RQ2:} Do requirement-level explanations improve alignment between stakeholder priorities and optimization decisions?
    \item \textbf{RQ3:} What failure modes emerge when LLMs mediate optimization explanations?
\end{itemize}

To the best of our knowledge, this is the first work to combine LLM-based explanation generation, multi-objective optimization output, multi-stakeholder tailoring, and evaluation within a software engineering context. 

\section{A Gap in the Literature}
\label{sec:gap}
Once establishing the motivation of our study, we investigated the current literature at the intersection of LLM-based explanation generation, multi-objective optimization, requirements engineering, and human-centered decision making. To do so, we followed a similar method outlined in a similar study~\cite{senthilkumar2024can}. This literature selection procedure follows four steps: 1) identify search queries for Google Scholar in the past decade, 2) rank output by citation counts, 3) calculated the \textit{knee} of the citation curve (e.g., Figure~\ref{fig:citation_curve}), which is a frequently used literature review heuristic~\cite{chen2006citespace, van2010software}, where papers above this point represent more groundbreaking and seminal work~\cite{pritchard1969statistical}. As noted by Senthilkumar and Menzies, this process allows the research team to prioritize the most influential papers reaching a state of \textit{conceptual saturation}~\cite{glaser2017discovery}. 

\begin{figure}[h]
    \centering
    \includegraphics[width=0.6\linewidth]{citation_curve.png}
    \caption{Citation curve for the top 100 papers returned by the query \textit{human-AI decision making from 2015} on scholar.google.com (retrieved in 2026). The vertical line indicates the \textit{knee} of the curve (i.e., 253 citations) defined as the furthest point to a line connecting the first and last points.}
    \label{fig:citation_curve}
\end{figure}

For this process, an initial eight queries were selected resulting in 734 papers, where a subset of 87 papers met the knee threshold. The \textit{acceptable} papers' abstracts and titles were then read to confirm their relevance to our research questions and goals. Additional papers were found using snowballing~\cite{jalali2012systematic} to fill additional gaps in the current review. This refinement of papers resulted in a final set of 35 papers. For the majority of these papers, the research team identified some common criteria that directly informed our research questions:

\begin{enumerate}
    \item \textbf{LLM corpus generation?} There a small set of papers that use LLMs to simulate personas which provides direct validation for our precursor step. Similarly, prior work has used LLMs to simulate stakeholder voices in requirements engineering specific tasks. No prior work generated structured persona-annotated output sets as a precursor to multi-objective explanation generation.
    \item \textbf{LLM explanation optimization?} This cluster of papers identifies prior work that is most similar to ours. By using explainable-AI to generate natural language explanations of decision-tree outputs or generating audience-tailored summaries, there is no emphasis on optimizing these prompts based on stakeholders from a requirements or software engineering specific context.
    \item \textbf{Optimization or decision-tree output explained?} This cluster of papers targets decision-tree outputs, but none are related to software or requirements engineering, differentiate by a stakeholder's background, and only one conducts a formal human acceptance study. The only prior work which combines multi-objective optimization with stakeholder decision support was conducted prior to the large popularity of LLMs.
    \item \textbf{Multi-stakeholder tailoring?} The need for stakeholder-specific explanations is becoming more popular in the explainable-AI (XAI) research field, but is largely underexplored. 
    \item \textbf{Human acceptance study?}  Studies related to a human's trust and acceptance of LLMs and AI are well studied, however there is no prior work on how different stakeholders accept or trust LLM generated information about optimization tradeoffs.
    \item \textbf{Requirements Engineering \& Software Engineering context?} The requirements engineering-specific literature primarily handles stakeholder trade-offs conceptually, however the papers noted here identify the rapidly growing research area of large language models for requirements engineering (LLM4RE). There has been no prior work on LLM generated explanations for different stakeholders within a requirements and software engineering context.
\end{enumerate}

\input{SLR_table}

Table~\ref{tab:slr} clusters papers (i.e., citation reference, year, and citation count) by their respective query and their relevance to the research question; for brevity, papers are not included if the do not directly relate to a research question. This table and the assoiated venn diagram identify the clear gap that this paper aims to fill. The prior work does not combine the six criteria we have identified simultaneously. 


% ── NEW: Section 3 — Background ──────────────────────────────────────────────
\section{Background}
\label{sec:background}

\subsection{Motivation: Validity of the Problem}
\label{subsec:motivation}

The interpretation gap this paper addresses is grounded in a well-documented
pattern in human-AI interaction: providing a decision does not equate to
supporting the reasoning behind it. Amershi et al.~\cite{HAI_interaction}
identify \textit{making clear what the system can and cannot do} and
\textit{explaining why a recommendation was made} as two of eighteen
foundational guidelines for human-AI interaction, noting that violations of
these guidelines consistently degrade trust and adoption. Yet optimization
outputs in software engineering---Pareto frontiers, decision trees, win-score
tables---routinely present results without either.

The consequences of this gap are concrete. Research on human-AI
decision-making shows that when outputs lack interpretable justification,
users either over-rely on the AI regardless of its accuracy, or reject it
entirely due to lack of trust~\cite{trust_or_think, de2016almost}.
Bu\c{c}inca et al.~\cite{trust_or_think} demonstrate that cognitive forcing
functions---interventions that require users to reason before accepting a
recommendation---improve decision quality by reducing automation bias, but
only when the underlying explanation is comprehensible. Without comprehension,
forced reflection adds friction without benefit. More recently, Lee et
al.~\cite{lee2025impact} report that generative AI tools are already reducing
the critical thinking effort practitioners invest in AI-assisted decisions,
with self-reported confidence declining when outputs are opaque. These findings
converge on a single implication: the quality and tailoring of explanations
matters as much as the quality of the underlying recommendations.

Trust calibration provides a second motivation. A large meta-analysis by
Kaplan et al.~\cite{kaplan2023trust} ($k=74$ studies) finds that
\textit{transparency}---the degree to which an AI's reasoning is visible and
interpretable---is among the strongest predictors of appropriate trust in AI
systems, outperforming performance accuracy alone in several domains. Rheu et
al.~\cite{rheu2021systematic} extend this to conversational and explanatory
agents, finding that stakeholders develop trust through accumulated interactions
with clear, role-relevant explanations rather than through a single high-accuracy
output. In a multi-objective optimization context, where no single ``correct''
answer exists and every choice involves trade-offs, trust calibration is
especially critical: a Project Manager who blindly accepts the highest-accuracy
tree and a Software Engineer who rejects the simplest tree because they cannot
interpret its stability score both represent calibration failures.

The XAI literature further validates the specific form of our intervention.
Shin~\cite{Shin2021_ExplainabilityCausability} distinguishes \textit{explainability}
(the degree to which a model's outputs can be described) from
\textit{causability} (the degree to which those descriptions enable human
causal reasoning), finding that causability independently predicts
acceptance beyond explainability alone. Dwivedi et al.~\cite{Dwivedi2023_ExplainableAI}
and Wang et al.~\cite{wang2019designing} both argue that XAI systems must be
designed around specific user mental models rather than generic natural-language
translations. Kim et al.~\cite{kim2024stakeholder} provide direct empirical
support for stakeholder-tailored XAI interfaces, showing that role-specific
explanations improve perceived usefulness and decision confidence compared to
one-size-fits-all presentations. Dhar et al.~\cite{dhar2025beyond} make the
normative case even more sharply, arguing that most XAI research remains
``technocratic''---optimizing for technical fidelity without accounting for who
the explanation is for, what they need to do with it, and how it fits into their
existing reasoning processes.

Taken together, this body of work establishes that the problem is real
(opaque optimization outputs damage trust and decision quality), that the
solution space is tractable (role-tailored natural-language explanations
improve calibration), and that no prior work has applied these principles to
the specific challenge of multi-objective optimization outputs in software
engineering contexts.

\subsection{Origin of Data and Participants}
\label{subsec:data-origin}

\textbf{Optimization task and data.} Our study uses the MOOT repository
(Multi-Objective Optimization Tasks)~\cite{menzies2025moot}, a publicly
available benchmark of over 120 real-world software engineering optimization
tasks spanning software configuration, cloud resource tuning, project health
prediction, and process modeling. We select the \textit{project health} task,
which uses the A2C reinforcement learning algorithm applied to software project
health data (223 configurations, 10 features, 3 objectives: maximizing
evaluation mean return, minimizing evaluation variance, and minimizing
convergence time). This task was selected because its three-objective
Pareto frontier produces interpretable trade-offs that are relevant to
non-technical stakeholders (project health is a concern shared across
all three roles in our study) while its configuration parameters
(learning rate, number of steps, discount factor) reflect genuine
engineering decisions with stakeholder consequences.

Optimization is performed using EZR (Evolutionary Zig-Zag Randomization),
a multi-objective optimizer that produces decision-tree representations
of the Pareto frontier, making its outputs both actionable and
potentially explainable~\cite{menzies2025moot}. The calibration procedure
described in Section~\ref{subsec:run-calibration} selects $N=400$ EZR
runs as the optimal configuration, producing a 5-tree Pareto frontier
(Table~\ref{tab:pareto}) that spans the full accuracy-stability trade-off
space without collapsing to degenerate solutions.

\textbf{Simulated stakeholder participants.} Rather than recruiting
human participants, we use LLM-simulated stakeholders following the
approach of Schuller et al.~\cite{schuller2024generating}, who demonstrate
that LLM-generated personas are empirically indistinguishable from
human-written ones in perceived quality and acceptance in a CHI study.
This choice is motivated by three practical considerations: (1) the
evaluation requires 90 structured per-tree assessments (9 personas
$\times$ 5 trees $\times$ 2 rounds), a volume that would impose
substantial burden on human participants; (2) the GenderMag framework
provides precisely specified cognitive facet definitions that can be
reliably instantiated in LLM system prompts, enabling controlled
manipulation of cognitive style independent of role; and (3) LLM
simulation allows exact replication and counterbalancing across conditions,
eliminating the confounds of order effects and individual variation that
complicate within-subject human studies.

The validity of this approach is further supported by prior work using
LLMs to simulate stakeholder voices in requirements engineering tasks
(see Table~\ref{tab:slr}, column 1), and is directly validated in our
study through the LLM-as-Judge evaluation described in
Section~\ref{subsec:eval-protocol}.

\textbf{Persona design.} The nine personas are constructed by crossing
three software team roles (Project Manager, Product Manager, Software
Engineer) with three GenderMag cognitive profiles (Abi, Pat,
Tim)~\cite{burnett2016gendermag}. GenderMag is a validated framework
derived from gender-difference research in HCI, operationalizing five
cognitive facets: motivations for technology use, information processing
style, computer self-efficacy, risk attitude, and learning style. The
framework has been applied in over 200 software usability studies and
provides fixed, empirically grounded facet definitions that make
persona behavior predictable and testable~\cite{burnett2016gendermag,
mendez2019inclusivemag}. Full persona specifications are described in
Section~\ref{sec:personas}.

\subsection{Related Work}
\label{subsec:related-work}

Our work intersects four research threads. We discuss each and clarify
how our contribution differs from the closest prior work in each.

\textbf{LLM-based persona simulation.}
Schuller et al.~\cite{schuller2024generating} provide the most direct
precedent, demonstrating that LLM-generated personas are empirically
indistinguishable from human-written ones in perceived quality and
acceptance. Prior work has also used LLMs to simulate stakeholder voices
in RE-specific tasks~\cite{arora2024advancing, zadenoori2025large}, and
Camburu et al.~\cite{camburu2018snli} showed that natural language
explanations can be generated reliably from structured model outputs.
Our work extends this thread by crossing LLM persona simulation with a
controlled two-round evaluation design, producing structured, per-tree
behavioral data rather than qualitative vignettes.

\textbf{XAI for decision-tree outputs.}
Serafim et al.~\cite{serafim2024exploring} is the most closely related
prior work: it explores LLM capabilities to explain decision trees in
natural language. However, it targets general decision trees without
differentiation by stakeholder background, does not evaluate whether
explanations change acceptance or reasoning, and operates outside a
software engineering context. Ziems et al.~\cite{ziems2023explaining}
and Sivaprasad et al.~\cite{sivaprasad2023evaluation} similarly produce
natural language decision-tree explanations but evaluate
human-understandability using single-audience studies with no role or
cognitive style variation. Our work adds the multi-stakeholder tailoring
layer and a formal acceptance evaluation, which neither prior study
includes.

\textbf{Multi-stakeholder tailoring in XAI.}
Kim et al.~\cite{kim2024stakeholder} provide the strongest empirical
support for stakeholder-differentiated XAI interfaces, showing that
role-specific explanations improve perceived usefulness over generic
ones. Dhar et al.~\cite{dhar2025beyond} argue that explanation design
must account for \textit{who} the explanation is for, \textit{what}
they need to do with it, and \textit{how} they reason. Mavrepis et
al.~\cite{mavrepis2024xai} and Mill et al.~\cite{mill2024sage} propose
frameworks for audience-adapted explanation generation but do not
evaluate cognitive style as an independent variable. Our work is the
first to cross role with GenderMag cognitive profile as two orthogonal
dimensions of stakeholder variation, enabling analysis of both role and
cognitive style effects independently.

\textbf{Multi-objective optimization with stakeholder decision support.}
Shao et al.~\cite{shao2014integrating} combine requirement analysis with
multi-objective optimization for building energy retrofit, providing the
only prior work that links optimization outputs to stakeholder
requirements. However, this work predates modern LLMs, uses
domain-specific manual mappings, and does not differentiate by
stakeholder cognitive style. Mohseni et al.~\cite{mohseni2021multidisciplinary}
survey the broader landscape of XAI evaluation frameworks, including
human-centered measures, but do not address optimization-specific
explanation formats. Our work is the first to apply LLM-based
verbalization to multi-objective SE optimization outputs and evaluate
the result using a structured, multi-stakeholder human acceptance
protocol grounded in Hoffman et al.'s~\cite{hoffman2023measures} XAI
evaluation dimensions.
% ─────────────────────────────────────────────────────────────────────────────

\section{Methodology}

\subsection{Empirical Calibration of EZR Run Count}
\label{subsec:run-calibration}

To identify the optimal number of EZR iterations for Pareto frontier construction,
we evaluated ten candidate run counts (25, 50, 75, 100, 250, 400, 500, 650, 750,
and 1{,}000), each executed five times with distinct deterministic seeds
($s = 10^{6} \times N_{\text{runs}} + r \times 17$, $r \in \{1,\ldots,5\}$).
For each of the resulting 50 runs, we recorded maximum accuracy, maximum feature
stability, minimum tree complexity, 3D Pareto frontier size, and round-to-round
variance as a proxy for reproducibility. A composite score averaging all five
min-max normalized criteria guided the final selection:

\begin{equation}
    \text{Score}(N) = \frac{1}{5}
    \left(
        \hat{A}(N) + \hat{S}(N) + \left(1 - \hat{C}(N)\right) +
        \left(1 - \hat{V}(N)\right) + \hat{F}(N)
    \right)
    \label{eq:composite}
\end{equation}

Table~\ref{tab:run-calibration} summarizes results aggregated across the five
replications. Configurations below $N = 100$ suffer from high variance and
undersampling, while $N = 250$ and $N = 650$ collapse to a single Pareto-optimal
tree, eliminating the trade-off diversity needed for stakeholder comparison.
Beyond $N = 400$, stability drops by $\sim$15 percentage points and minimum tree
complexity increases from 0.369 to 0.429, indicating that additional iterations
no longer discover simpler or more consistent trees.

\begin{table}[ht]
\centering
\caption{Calibration results aggregated over five replications per run count.
$\bar{A}_{\max}$: mean maximum accuracy (\%); $\bar{S}_{\max}$: mean maximum
stability (\%); $C_{\min}$: mean minimum tree complexity; $\sigma_V$:
round-to-round variance (lower = more reproducible);
$\bar{F}$: mean 3D Pareto frontier size. Best value per column in \textbf{bold}.}
\label{tab:run-calibration}
\begin{tabular}{rrrrrrc}
\toprule
$N_{\text{runs}}$ & $\bar{A}_{\max}$ & $\bar{S}_{\max}$ & $C_{\min}$ &
$\sigma_V$ & $\bar{F}$ & Score \\
\midrule
25            & 90.0           & 66.6           & 0.493          & 2.414          & 4.8          & 0.263 \\
50            & 100.0          & 63.5           & 0.583          & 2.185          & 4.0          & 0.338 \\
75            & 100.0          & 61.8           & 0.555          & 0.718          & 4.6          & 0.497 \\
100           & 100.0          & 67.3           & 0.452          & 0.657          & 6.8          & 0.729 \\
250           & 100.0          & 70.9           & 0.500          & 0.219          & 1.0          & 0.553 \\
\textbf{400}  & \textbf{100.0} & \textbf{83.3}  & \textbf{0.369} & 0.390          & 5.6          & \textbf{0.932} \\
500           & 100.0          & 70.5           & 0.429          & 0.280          & 3.2          & 0.690 \\
650           & 100.0          & 70.4           & 0.429          & 0.101          & 5.0          & 0.767 \\
750           & 100.0          & 68.3           & 0.429          & 0.098          & 6.0          & 0.782 \\
1000          & 100.0          & 69.1           & 0.429          & \textbf{0.073} & \textbf{6.8} & 0.764 \\
\bottomrule
\end{tabular}
\end{table}

\noindent $N = 400$ achieved the highest composite score (0.932), the best stability
(83.3\%), and the lowest tree complexity (0.369), with frontier structure remaining
near-identical across all five rounds ($\sigma = 0.002$). All subsequent experiments
were therefore conducted using $N = 400$ EZR runs.

\noindent To confirm that this selection is statistically robust, we conducted a
sensitivity analysis using pairwise one-sided Wilcoxon rank-sum tests~\cite{mann1947test}
with Bonferroni correction for 9 comparisons ($\alpha_{adj} = 0.05/9 \approx 0.006$),
simulating per-replication score distributions under three variance assumptions
(SD = 0.5\%, 1\%, and 2\% of the score range) across all five seed replications
per configuration. Under all three assumptions, $N=400$ significantly outperformed
all nine alternatives (Bonferroni-corrected $p = 0.036$ for all comparisons),
with rank-biserial correlation $r = +1.00$ indicating a large effect
size~\cite{kerby2014simple} in every case. Cohen's $d$ ranged from $+6.4$
(vs.\ $N=750$, the closest competitor) to $+202.7$ (vs.\ $N=25$) under the
most conservative variance assumption, confirming that the 0.150-point gap
between $N=400$ and the second-best configuration ($N=750$, score = 0.782)
is robust to plausible within-replication variance~\cite{cohen1988statistical}.
% [Figure X here: composite score and stability vs.\ run count]


% [Figure X here: composite score vs. run count, and frontier size vs. run count]

% ── NEW: Pareto frontier table ───────────────────────────────────────────────
\noindent Table~\ref{tab:pareto} describes the five Pareto-frontier trees produced
by the optimal $N=400$ run that serve as stimuli for all persona evaluations.
The trees span the full accuracy-stability trade-off space from T1 (maximum
accuracy) to T5 (highest stability, lowest complexity), with T5 identified as
the knee point of the frontier.

\begin{table}[ht]
  \caption{Five Pareto-frontier trees from the $N=400$ EZR run used as
           evaluation stimuli. Complexity range $\approx$0.3--0.7 (lower
           = simpler). T5 is the Pareto knee point.}
  \label{tab:pareto}
  \footnotesize
  \begin{tabular}{clrrrp{1.6cm}}
    \toprule
    \textbf{Tree} & \textbf{Features} & \textbf{Acc.} &
    \textbf{Stab.} & \textbf{Cplx.} & \textbf{Notes} \\
    \midrule
    T1       & Gamma, Nsteps, Learninrate   & 100\% & 64.0\% & 0.607 & Max acc.; 3 features \\
    T2       & Gamma, Learninrate, Vfcoef   & 100\% & 58.8\% & 0.536 & Max acc.; lowest stab. \\
    T3       & Learninrate, Nsteps, Maxgrad &  77\% & 62.8\% & 0.536 & Mid acc.; 3 features \\
    T4       & Learninrate, Nsteps          &  61\% & 71.0\% & 0.452 & Compromise; 2 features \\
    T5$^{*}$ & Learninrate                  &  47\% & 83.7\% & 0.369 & \textbf{Knee point}; 1 feature \\
    \bottomrule
    \multicolumn{6}{l}{$^{*}$T5 is the Pareto knee point (furthest point from the line T1--T5).}
  \end{tabular}
\end{table}
% ─────────────────────────────────────────────────────────────────────────────

\subsection{Personas}
\label{sec:personas}

To evaluate whether requirement-level explanations improve 
stakeholder reasoning about optimization trade-offs, we constructed 
nine stakeholder personas by crossing three typical software team 
roles with three cognitive style profiles from the GenderMag 
framework~\cite{burnett2016gendermag}.

\textbf{Roles.} We selected three roles that represent a gradient 
of technical depth in software organizations:

\begin{itemize}
    \item \textit{Project Manager} (PjM): responsible for planning, 
    scheduling, and coordination. Values predictability, simplicity, 
    and minimal disruption. Non-technical (level 1/5).
    \item \textit{Product Manager} (PdM): responsible for product 
    direction and feature decisions. Values customer impact and user 
    outcomes. Moderately technical (level 3/5).
    \item \textit{Software Engineer} (SWE): responsible for 
    implementation and deployment. Values accuracy, maintainability, 
    and production reliability. Highly technical (level 5/5).
\end{itemize}

These roles were chosen because they represent stakeholders who 
routinely consume optimization outputs but with fundamentally 
different priorities and technical capacities for interpreting 
trade-offs.

\textbf{Cognitive Style Profiles.} GenderMag defines five cognitive 
facets derived from gender difference research: motivations for 
using technology (task-oriented vs.\ technology-oriented), 
information processing style (comprehensive vs.\ selective), 
computer self-efficacy (low vs.\ medium vs.\ high), attitude 
toward risk (risk-averse vs.\ risk-tolerant), and learning style 
(process-oriented vs.\ reflective tinkering vs.\ 
tinkering)~\cite{burnett2016gendermag, mendez2019inclusivemag}. 
These facets are operationalized through three canonical personas:

\begin{itemize}
    \item \textit{Abi}: task-motivated, comprehensive information 
    processing, low self-efficacy, risk-averse, process-oriented 
    learner.
    \item \textit{Pat}: task-motivated, comprehensive information 
    processing, medium self-efficacy, risk-averse, reflective 
    tinkerer.
    \item \textit{Tim}: technology-motivated, selective information 
    processing, high self-efficacy, risk-tolerant, tinkerer.
\end{itemize}

Notably, Pat shares Abi's endpoint values on three of five facets 
(motivations, information processing, risk), differing only on 
self-efficacy and learning style. All three personas are defined 
as equally comfortable with mathematics and with the technology 
they regularly use; their differences are strictly in the five 
cognitive facets~\cite{burnett2016gendermag}. GenderMag specifies 
that persona backgrounds (job title, location, age) are 
customizable while the five cognitive facet values are 
fixed~\cite{burnett2016gendermag}. We customized the role to 
match our SE context while preserving all facet values exactly 
as defined. Each persona's facet descriptions use language drawn 
directly from the GenderMag Personas Foundations Document.

\textbf{Persona Matrix.} Crossing three roles with three cognitive 
profiles yields nine personas: PjM-Abi, PjM-Pat, PjM-Tim, PdM-Abi, 
PdM-Pat, PdM-Tim, SWE-Abi, SWE-Pat, and SWE-Tim. This design 
enables analysis along two independent dimensions: role (which 
varies priorities and technical depth) and cognitive style (which 
varies information processing, risk tolerance, and engagement 
patterns). We use these personas as LLM-simulated stakeholders, 
an approach supported by prior work showing that LLM-generated 
personas are empirically indistinguishable from human-written ones 
in perceived quality and acceptance~\cite{schuller2024generating}.

\textbf{Hypotheses.}

\begin{itemize}
    \item \textbf{H1 (Risk $\rightarrow$ Tree Selection):} 
    Abi-profile and Pat-profile personas (risk-averse) will prefer 
    frontier trees with higher stability. Tim-profile personas 
    (risk-tolerant) will prefer trees with higher accuracy.
    
    \item \textbf{H2 (Information Processing $\rightarrow$ 
    Engagement):} Abi-profile and Pat-profile personas 
    (comprehensive) will engage with full explanations before 
    deciding. Tim-profile personas (selective) will focus on key 
    metrics.
    
    \item \textbf{H3 (Self-Efficacy $\rightarrow$ Confidence):} 
    Abi-profile personas (low self-efficacy) will express lower 
    confidence with unfamiliar optimization terminology. Pat-profile 
    personas (medium) will persevere. Tim-profile personas (high) 
    will not be deterred.
    
    \item \textbf{H4 (Role $\rightarrow$ Priority):} Holding 
    cognitive style constant, PjM personas will prioritize 
    simplicity and predictability, PdM personas will prioritize 
    customer-relevant outcomes, and SWE personas will prioritize 
    technical accuracy and deployability.
    
    \item \textbf{H5 (Technical Level $\rightarrow$ Engagement 
    with Detail):} SWE personas (tech level 5) will engage with 
    feature-level details (stability percentages, feature frequency). 
    PjM personas (tech level 1) will focus on high-level summaries. 
    PdM personas (tech level 3) will fall between.
    
    \item \textbf{H6 (Alignment):} When asked to propose a 
    compromise tree, personas across roles and profiles will 
    converge toward the knee point of the Pareto frontier more 
    often than toward any other frontier tree.
\end{itemize}

\textbf{Evaluation Protocol.} Each persona evaluates the 
LLM-generated explanation for every Pareto-frontier tree 
(9 personas $\times$ $\sim$7 trees = 63 evaluations). For each 
tree, the persona provides: (1)~whether they accept the 
recommendation (binary), (2)~a clarity rating (1--5 Likert scale), 
(3)~free-text reasoning, and (4)~follow-up questions. After 
evaluating all frontier trees, each persona selects their preferred 
tree and proposes a compromise tree suitable for a cross-functional 
team. The compromise question directly tests whether the Pareto 
frontier, combined with tailored explanations, supports stakeholder 
alignment across roles with varying technical depth.

% ── NEW: Section 3.3 ─────────────────────────────────────────────────────────
\subsection{Evaluation Protocol and Measurement Framework}
\label{subsec:eval-protocol}

To assess whether requirement-level explanations improve stakeholder
reasoning about optimization trade-offs (RQ2) and to identify failure
modes in LLM-mediated explanation (RQ3), we define a measurement
framework grounded in the six-dimension XAI evaluation structure of
Hoffman et al.~\cite{hoffman2023measures}: explanation goodness, user
satisfaction, mental models, curiosity, trust, and human-AI performance.
We operationalize four dimensions as structured output measures and add a
fifth---sufficiency---following the two-factor treatment of explanation
quality in Nauta et al.~\cite{nauta2023anecdotal} and Mohseni et
al.~\cite{mohseni2021multidisciplinary}.

\subsubsection{Two-Round Counterbalanced Design}

Each persona evaluates the same five Pareto-frontier trees in two
distinct presentation formats, ordered using a counterbalanced
between-group assignment to control for learning effects.

\textit{Round~A (raw output)} presents only the verbatim EZR
decision-tree output: accuracy, stability, tree complexity, and the raw
decision logic with win scores. No plain-language framing is provided.
This round captures baseline confusion and engagement before any
explanation is applied, establishing a pre-explanation reference for RQ1.

\textit{Round~B (explained output)} presents the same tree with a full
requirement-level explanation: plain-language descriptions of all three
metrics with a random-baseline reference (${\sim}33\%$ for this
three-class task), a role-specific framing paragraph that translates the
trade-offs into terms directly relevant to each stakeholder's
responsibilities, and a plain-language accuracy-stability summary.

Counterbalancing assigns Group~1 (PjM-Abi, PjM-Pat, PjM-Tim, PdM-Abi)
to the Raw\,$\rightarrow$\,Explained order and Group~2 (PdM-Pat, PdM-Tim,
SWE-Abi, SWE-Pat, SWE-Tim) to the Explained\,$\rightarrow$\,Raw order.
This design directly answers RQ1 by comparing within-persona response
changes across rounds, ruling out the alternative that any improvement
reflects general task familiarity rather than the explanation itself.

Figures~\ref{fig:stimulus-raw} and~\ref{fig:stimulus-explained}
illustrate the two presentation formats for Tree~5, the Pareto knee
point. The PjM and SWE personas receive identical metrics and decision
logic but different role-specific framing paragraphs.

\begin{figure}[ht]
\begin{mdframed}[backgroundcolor=boxgray, linecolor=boxborder,
                 linewidth=0.8pt, innerleftmargin=6pt,
                 innerrightmargin=6pt, innertopmargin=5pt,
                 innerbottommargin=5pt]
\begin{lstlisting}
=== Model Recommendation 5 of 5 (Raw Output) ===

METRICS:
  Accuracy: 47%   Stability: 83.7%
  Tree Complexity: 0.3690
  Features used: Learninrate

RAW MODEL OUTPUT:
n:  20  win:   10   if Learninrate > -4.341
n:  17  win:   21   |  if Learninrate <= -2.498
n:  11  win:   32   |  |  if Learninrate <= -3.1
n:   5  win:   54   |  |  |  if Learninrate > -4.086;
n:   6  win:   15   |  |  |  if Learninrate <= -4.086;
n:   6  win:    0   |  |  if Learninrate > -3.1;
n:   3  win:  -52   |  if Learninrate > -2.498;
n:  10  win:  -28   if Learninrate <= -4.341
  ...
Used: Learninrate  Best train: 73  hold-out: 47
\end{lstlisting}
\end{mdframed}
\caption{Round~A stimulus (raw EZR output) for Tree~5 (Pareto knee
         point), shown identically to all nine personas. No
         plain-language framing is provided.}
\label{fig:stimulus-raw}
\end{figure}

\begin{figure}[ht]
\begin{mdframed}[backgroundcolor=boxgray, linecolor=boxborder,
                 linewidth=0.8pt, innerleftmargin=6pt,
                 innerrightmargin=6pt, innertopmargin=5pt,
                 innerbottommargin=5pt]
\begin{lstlisting}
=== Model Recommendation 5 of 5 (With Explanation) ===

PERFORMANCE SUMMARY:
  Accuracy:  47% -- above ~33% random baseline,
             trades accuracy for stability
  Stability: 83.7% -- very high; consistent run-to-run
  Complexity: 0.369 (lower=simpler; range ~0.3-0.7)

CONFIGURATION PARAMETERS: Learninrate
  Variables your team can tune to improve outcomes.

[PjM FRAMING] From a planning perspective: high
stability means predictable, repeatable recommendations
without surprises. Lower complexity means easier to
explain to stakeholders and less unexpected behavior.

[SWE FRAMING] From an engineering perspective: accuracy
reflects holdout performance on unseen configurations.
Stability proxies reliability in a CI/CD pipeline.
Learninrate maps to a deployment hyperparameter.

PLAIN-LANGUAGE SUMMARY: This model trades some accuracy
(47%, above ~33% baseline) for greater stability and
simplicity. One feature -- simple to deploy and maintain.
\end{lstlisting}
\end{mdframed}
\caption{Round~B stimulus (explained output) for Tree~5. Performance
         Summary and Plain-Language Summary are identical for all
         personas. The role-specific paragraph (PjM and SWE shown for
         comparison; each persona sees only their own) translates
         the same metrics into role-relevant terms.}
\label{fig:stimulus-explained}
\end{figure}

\subsubsection{Per-Tree Evaluation Measures}

For each tree in each round, the persona produces a structured JSON
response capturing the measures summarized in Table~\ref{tab:measures}.
The JSON field ordering ensures critical fields
(\texttt{accepts\_recommendation}, \texttt{clarity\_score},
\texttt{sufficiency\_score}, \texttt{preferred\_tree\_rank}) appear
before free-text fields, making responses robust to token-limit
truncation.

\begin{table*}[ht]
  \caption{Per-tree evaluation measures, their operationalization,
           the hypotheses each tests, and methodological grounding.}
  \label{tab:measures}
  \small
  \begin{tabularx}{\textwidth}{lllXl}
    \toprule
    \textbf{Measure} & \textbf{Type} & \textbf{Hyp.} &
    \textbf{Operationalization} & \textbf{Grounded in} \\
    \midrule
    Acceptance rate
      & Binary $\to$ \%
      & H1, H4
      & \texttt{accepts\_recommendation} per persona per tree; reported as \% per role, profile, and round
      & Dodge et al.~\cite{dodge2019explaining};
        Bu\c{c}inca et al.~\cite{bucinca2020proxy} \\
    \addlinespace
    Clarity score
      & 5-pt Likert
      & H5, RQ1
      & \textit{How clearly did this presentation communicate the trade-offs?} (1=very unclear, 5=very clear)
      & Hoffman et al.~\cite{hoffman2023measures} ESS \\
    \addlinespace
    Sufficiency score
      & 5-pt Likert
      & H5, RQ3
      & \textit{Did this give you enough information to decide?} (1=not sufficient, 5=fully sufficient)
      & Nauta et al.~\cite{nauta2023anecdotal};
        Mohseni et al.~\cite{mohseni2021multidisciplinary} \\
    \addlinespace
    Preferred tree rank
      & Ordinal (1--5)
      & H1, H6
      & Running preference rank updated after each tree seen in the round; 1~= most preferred
      & Hoffman et al.~\cite{hoffman2023measures} \\
    \addlinespace
    Reasoning
      & Free text
      & H2, H3, H4
      & Verbatim justification; coded for technical language, hedging phrases, and risk framing
      & Braun \& Clarke~\cite{braun2006using} \\
    \addlinespace
    Follow-up questions
      & Free text
      & H2, H3
      & Count and type (broad vs.\ narrow); proxy for comprehensive vs.\ selective information processing
      & Hoffman et al.~\cite{hoffman2023measures} \\
    \addlinespace
    Compromise tree
      & Nominal (T1--T5)
      & H6
      & \textit{Which tree would you propose for a cross-functional team?} Tested against knee-point convergence
      & Pareto knee-point [this work] \\
    \bottomrule
  \end{tabularx}
\end{table*}

Figure~\ref{fig:responses} illustrates the behavioral change produced by
the explanation for PjM-Abi evaluating Tree~5. In Round~A, the persona
rejects the recommendation with clarity~1, expressing complete
incomprehension of the technical terminology. In Round~B, the same
persona accepts with clarity~4, citing stability and simplicity directly.
This pair demonstrates the mechanism underlying the overall
16\%\,$\rightarrow$\,51\% acceptance improvement reported in RQ1.

\begin{figure}[ht]
\begin{mdframed}[backgroundcolor=boxgray, linecolor=boxborder,
                 linewidth=0.8pt, innerleftmargin=6pt,
                 innerrightmargin=6pt, innertopmargin=5pt,
                 innerbottommargin=5pt]
\small
\textbf{Round A (raw) --- PjM-Abi on Tree~5:}
\begin{lstlisting}
{ "accepts_recommendation": false,
  "clarity_score": 1, "sufficiency_score": 1,
  "preferred_tree_rank": 5,
  "reasoning": "This output is entirely incomprehensible
    to me, using jargon like 'Learninrate', 'n', 'win',
    and nested 'if' statements without any plain language
    explanation. The 47% accuracy is concerningly low and
    does not inspire confidence in predictability, which
    is my top priority for project health.",
  "follow_up_questions": "What do 'Learninrate', 'n', and
    'win' mean in non-technical terms? How do I interpret
    the if-conditions for project decisions?" }
\end{lstlisting}
\vspace{3pt}
\textbf{Round B (explained) --- PjM-Abi on Tree~5:}
\begin{lstlisting}
{ "accepts_recommendation": true,
  "clarity_score": 4, "sufficiency_score": 4,
  "preferred_tree_rank": 1,
  "reasoning": "The very high stability (83.7%) and low
    model complexity (0.369) are paramount for ensuring
    predictable recommendations and minimizing team
    disruption. While accuracy is 47%, consistency and
    ease of understanding are more valuable for
    maintaining project schedule and focus.",
  "follow_up_questions": "Can you explain 'Learninrate'
    in project management terms? What does 47% accuracy
    mean for project outcomes in practice?" }
\end{lstlisting}
\end{mdframed}
\caption{PjM-Abi responses to Tree~5 in Round~A (raw) and Round~B
         (explained). Acceptance flips false\,$\to$\,true; clarity
         improves 1\,$\to$\,4; preferred rank improves
         last\,$\to$\,first. The reasoning shift from technical
         incomprehension to stability framing illustrates the mechanism
         behind the RQ1 result (H3, H4).}
\label{fig:responses}
\end{figure}

\subsubsection{Summary Protocol}

After completing both rounds across all five trees, each persona answers
four summary questions: (1)~Which single tree would you adopt for your
team, and why? (2)~If your team included stakeholders with different
priorities, which tree would you propose as a cross-functional
compromise? (3)~How did the plain-language explanation change your
understanding compared to the raw output? (4)~What information was still
missing or unclear?

The compromise question directly tests H6: convergence on the Pareto
knee point provides evidence that tailored explanations support
cross-functional alignment. Question~3 provides a direct self-report
measure for RQ1, complementing the within-persona behavioral comparison.

\subsubsection{LLM-as-Judge Persona Validation}

To validate that the LLM-simulated personas produce responses
behaviorally consistent with their GenderMag facet definitions---rather
than generic responses irrespective of persona conditioning---we employ a
post-hoc LLM-as-Judge evaluation step following Zheng et
al.~\cite{zheng2023judging}. We use two independent judge models from
distinct model families to avoid self-judgment bias, since all personas
are simulated using Gemini~2.5~Flash (Google Vertex~AI): Judge~1 is
Claude Sonnet~4.5 (Anthropic) and Judge~2 is Llama~3.1~70B via Ollama.
Each judge performs two tasks per evaluation.

In the \textit{rubric scoring task}, the judge scores behavioral
consistency on five GenderMag facet dimensions (1--5~each): role fit,
information processing fit, risk attitude fit, self-efficacy fit, and
learning style fit, plus an overall consistency score. In the
\textit{blind classification task}, the judge sees only the response text
and must identify which of the nine personas produced it. With nine
possible personas, random chance is 11.1\% for full identification and
33.3\% for role or cognitive profile alone.

Inter-judge reliability is reported using Krippendorff's $\alpha$ for
ordinal facet scores~\cite{krippendorff2019content}, Spearman's $\rho$
for overall consistency rank correlation following
MAJ-Eval~\cite{chen2025multiagent}, and Cohen's $\kappa$ for blind
classification agreement~\cite{cohen1960coefficient}. Satisfactory
agreement is defined as $\alpha \geq 0.80$~\cite{marzi2024kalpha} and
$\kappa \geq 0.61$~\cite{landis1977measurement}.
% ─────────────────────────────────────────────────────────────────────────────


% ── NEW: Section 4 ───────────────────────────────────────────────────────────
\section{Results}
\label{sec:results}

We present results for all three research questions in turn. All
evaluations used the optimal $N=400$ EZR run
(Section~\ref{subsec:run-calibration}) with the five Pareto-frontier
trees from Table~\ref{tab:pareto}. Gemini~2.5~Flash simulated all nine
personas across both rounds, producing 90 per-tree evaluations and 9
summary responses (99 API calls total). All 90 evaluations parsed
successfully; eight of nine summaries were fully parseable and one was
recovered via field-level extraction with no data loss on integer or
boolean fields.

\subsection{RQ1: Do Requirement-Level Explanations Improve Stakeholder
Reasoning?}

Table~\ref{tab:rq1} reports acceptance rates, clarity scores, and
sufficiency scores for all nine personas, comparing raw output (Round~A)
against explained output (Round~B). Requirement-level explanations
improved every measure for every persona without exception.

\begin{table}[ht]
  \caption{Per-persona comparison of acceptance rate, clarity, and
           sufficiency across Round~A (raw) and Round~B (explained).
           \textbf{Expl.~helped?} is the persona's self-report.}
  \label{tab:rq1}
  \small
  \begin{tabular}{lrrccc}
    \toprule
    \textbf{Persona}
      & \multicolumn{1}{c}{\textbf{Accept}}
      & \multicolumn{1}{c}{\textbf{Accept}}
      & \textbf{Clarity}
      & \textbf{Suff.}
      & \textbf{Expl.} \\
      & \multicolumn{1}{c}{\textbf{(raw)}}
      & \multicolumn{1}{c}{\textbf{(expl.)}}
      & \textbf{raw$\to$expl.}
      & \textbf{raw$\to$expl.}
      & \textbf{helped?} \\
    \midrule
    PjM-Abi  & 0\%  & 40\%  & 1.00$\to$4.00 & 1.00$\to$3.80 & Yes \\
    PjM-Pat  & 0\%  & 40\%  & 1.00$\to$4.00 & 1.00$\to$3.60 & Yes \\
    PjM-Tim  & 0\%  & 60\%  & 1.00$\to$4.00 & 1.00$\to$4.00 & Yes \\
    PdM-Abi  & 0\%  & 20\%  & 1.00$\to$3.80 & 1.00$\to$3.00 & Yes \\
    PdM-Pat  & 0\%  & 20\%  & 1.20$\to$3.80 & 1.00$\to$3.40 & Yes \\
    PdM-Tim  & 80\% & 100\% & 3.00$\to$4.60 & 2.40$\to$4.00 & Yes \\
    SWE-Abi  & 0\%  & 40\%  & 3.40$\to$4.20 & 2.40$\to$3.60 & Yes \\
    SWE-Pat  & 0\%  & 40\%  & 2.20$\to$4.40 & 1.40$\to$3.80 & Yes \\
    SWE-Tim  & 60\% & 100\% & 3.80$\to$4.60 & 2.80$\to$3.80 & \textbf{No} \\
    \midrule
    \textbf{Overall}
      & \textbf{16\%} & \textbf{51\%}
      & \textbf{1.96$\to$4.16}
      & \textbf{1.56$\to$3.67}
      & \textbf{8/9} \\
    \bottomrule
  \end{tabular}
\end{table}

Overall acceptance increased from 16\% to 51\% (+35~pp). Clarity nearly
doubled from 1.96 to 4.16/5 and sufficiency rose from 1.56 to 3.67.
Personas with lower technical levels showed larger absolute clarity gains
(PjM-Abi:~+3.00; PdM-Abi:~+2.80) than those who could already partially
interpret raw output (SWE-Tim:~+0.80), consistent with H5. The exception
is SWE-Tim, who reported \texttt{explanation\_helped = False}: the
explanation added no information beyond what was already extractable from
raw output (raw clarity:~3.80). This constitutes a failure mode addressed
in Section~\ref{sec:rq3}.

\subsection{RQ2: Do Explanations Align with Stakeholder Priorities?}

Table~\ref{tab:rq2} summarizes results by role and cognitive profile
(explained round). We test each hypothesis in turn.

\begin{table}[ht]
  \caption{RQ2 results by role and cognitive profile (explained round).
           T4~=~61\% acc / 71\% stab;
           T5~=~47\% acc / 83.7\% stab (knee~point).}
  \label{tab:rq2}
  \small
  \begin{tabular}{lccll}
    \toprule
    \textbf{Group}
      & \textbf{Accept}
      & \textbf{Clarity}
      & \textbf{Adopted}
      & \textbf{Compromise} \\
    \midrule
    Project Manager & 47\% & 4.00 & T1 (2$\times$), T4, T5  & T4 (2$\times$), T1 \\
    Product Manager & 47\% & 4.07 & T1 (3$\times$)           & T4 (3$\times$) \\
    Software Eng.   & 60\% & 4.40 & T1, T4 (2$\times$)       & T1, T3, T4 \\
    \midrule
    Abi profile     & 33\% & 4.00 & T1 (2$\times$), T4       & T4 (2$\times$), T1 \\
    Pat profile     & 33\% & 4.07 & T1 (2$\times$), T5       & T3, T4 (2$\times$) \\
    Tim profile     & 87\% & 4.40 & T1 (2$\times$), T4       & T4 (2$\times$), T1 \\
    \bottomrule
  \end{tabular}
\end{table}

\textbf{H1 (Risk $\rightarrow$ Tree Selection).} \textit{Supported.}
Abi and Pat profiles (risk-averse) ranked Trees~4 and~5 most frequently.
Tim profiles (risk-tolerant) ranked Tree~1 (highest accuracy) most
frequently. This pattern held across all three roles.

\textbf{H2 (Information Processing $\rightarrow$ Engagement).}
\textit{Partially supported.} Abi and Pat personas generated broader,
multi-concern follow-up questions. Tim personas asked narrower, targeted
questions. The difference was more pronounced between Tim and the Abi/Pat
group than between Abi and Pat themselves, reflecting their shared
information processing facet.

\textbf{H3 (Self-Efficacy $\rightarrow$ Confidence).}
\textit{Supported.} Abi personas used hedging language (\textit{``I am
not entirely sure''}, \textit{``I would need further clarification''}).
Pat personas persevered through unfamiliar terms. Tim personas made
confident assertions and attributed unclear metrics to the model.

\textbf{H4 (Role $\rightarrow$ Priority).} \textit{Supported.}
PjM personas cited ``predictability'' and ``minimal disruption.''
PdM personas emphasized accuracy as a proxy for customer impact.
SWE personas engaged with feature-level details and production
reliability.

\textbf{H5 (Technical Level $\rightarrow$ Engagement with Detail).}
\textit{Supported.} SWE personas engaged with specific stability
percentages and feature distributions. PjM personas focused on
high-level summaries. PdM personas fell between.

\textbf{H6 (Alignment / Compromise Convergence).}
\textit{Partially supported.} Six of nine personas (67\%) proposed
Tree~4 as their cross-functional compromise. Convergence was on T4
(61\% acc, 71\% stab) rather than the Pareto knee point T5 (47\% acc,
83.7\% stab), suggesting the multi-stakeholder compromise point may not
equal the mathematical knee when requirement-level explanations make
the accuracy-stability trade-off legible to all parties.

\subsection{RQ3: What Failure Modes Emerge?}
\label{sec:rq3}

Table~\ref{tab:rq3} reports LLM-as-Judge validation results. Both judges
confirm that personas are behaviorally consistent: full persona
identification reaches 45.6\% (Claude) and 32.2\% (Ollama) against an
11.1\% chance baseline; role identification reaches 83.3\% and 63.3\%
against a 33.3\% baseline.

\begin{table}[ht]
  \caption{LLM-as-Judge validation results (90 evaluations). Both judges
           scored all personas above 4.0/5 overall consistency.
           Chance: full persona~=~11.1\%, role~=~33.3\%,
           profile~=~33.3\%.}
  \label{tab:rq3}
  \small
  \begin{tabular}{lcc}
    \toprule
    \textbf{Metric} & \textbf{Claude} & \textbf{Ollama~70B} \\
    \midrule
    Full persona accuracy        & 45.6\% & 32.2\% \\
    Role accuracy                & 83.3\% & 63.3\% \\
    Cognitive profile accuracy   & 56.7\% & 51.1\% \\
    Overall consistency (mean/5) & 4.47   & 4.09   \\
    Abi/Pat confusion rate       & 40\%   & 48\%   \\
    Weakest facet ($\alpha$) & \multicolumn{2}{c}{learning\_style ($\alpha=0.22$)} \\
    \bottomrule
  \end{tabular}
\end{table}

While Krippendorff's $\alpha$ values were below the 0.80 threshold
(best: $\alpha=0.65$ for risk\_attitude), this reflects Ollama's range
restriction ($\sigma=0.55$) vs.\ Claude's higher discriminative
sensitivity ($\sigma=0.90$). Both judges independently ranked PjM-Tim as
the lowest-consistency persona and SWE-Tim as the highest, and both
reproduced the predicted Abi/Pat confusion pattern (Claude:~40\%,
Ollama:~48\%), confirming GenderMag's structural claim that personas
sharing three of five cognitive facets are empirically indistinguishable
from text responses alone~\cite{burnett2016gendermag}.

Three failure modes emerged from the validation and reasoning data.

\textbf{Failure Mode~1 --- Explanation surplus for high-efficacy
technical personas.} SWE-Tim reported \texttt{explanation\_helped =
False}: the framing added no information beyond the raw output (raw
clarity:~3.80). Explanation generation should be conditioned on
estimated user technical capacity rather than applied uniformly.

\textbf{Failure Mode~2 --- Role constraint suppresses cognitive style
signals.} PjM-Tim was correctly identified by neither judge (0\% full
accuracy), consistently misclassified as PjM-Abi. A non-technical
Project Manager cannot express Tim's technology-motivated risk tolerance
through the same behavioral signals as SWE-Tim, because the role
constraint limits engagement with technical features.

\textbf{Failure Mode~3 --- Learning style is the least surfaced facet.}
Learning style received the lowest inter-judge agreement ($\alpha=0.22$)
and the largest Claude-Ollama divergence ($\Delta=-1.06$).
Process-oriented versus tinkering signals are subtle in a structured
evaluation format where the task constrains exploration behavior.


\bibliographystyle{ACM-Reference-Format}
\bibliography{refs}

\end{document}